\section{Introduction}

Set-asides open SME access to public contracts by closing the auction to
everyone else. The design is mechanical: protected firms come in, unprotected
firms go out. Removing rival bidders weakens the competitive discipline that
forms the price. Admitting more SMEs into a now-restricted contest pushes
back. The welfare case depends entirely on which force is larger, and on how
the protected SME pool reorganizes once the auction is closed to non-SMEs.
The empirical literature has documented the cost of SME-favoring rules but
rarely separates these two channels, even though policy design lives in their
balance.

I show that a set-aside replaces competition with eligibility, and that when
the bidders it removes are the price-forming ones the replacement is partial,
expensive, and dominated by a price preference. The point is not that SME
set-asides can be costly: \citet{marion2007} and \citet{nakabayashi2013}
document that they often are. \citet{athey2013} and
\citet{krasnokutskayaseim2011} go further and show that the design of SME
support shapes whether the cost falls on allocation, entry, or both. None of
these decomposes the price effect into what an exclusionary rule does to the
price-forming pool and what the protected pool does in response. Without that
split, the policy debate runs as SME support versus no SME support. The
relevant comparison is different: support that preserves the price-forming
pool against support that achieves access by deleting rival bidders. This
paper places a real procurement reform inside that comparison.

The setting is S\~ao Paulo's centralized electronic procurement platform, the
Bolsa Eletr\^onica de Compras (BEC). In late 2017 and early 2018, a legal
reinterpretation moved the state's SME-only eligibility threshold from the
total value of a purchase notice to the item level. The change mechanically
pushed medical and hospital supplies (Group 65) into SME-only tendering,
because such purchases combine many small items inside large notices. The
empirical cutoff is \policyCutoffMonth{}, when Group-65 public buyers began
mass adoption of the new functionality. The trigger was legal and
administrative, not a medical-supply-specific policy, and it changed
admissibility rather than the production technology of the goods.

The first-stage facts fit the institutional account. After the cutoff,
SME-only adoption in Group 65 rises sharply, non-SME participation falls, SME
participation roughly doubles, and winning prices rise. A compact
difference-in-differences benchmark against \controlGroupCount{} never-treated
product groups places the associated price movement at
\didLowerPct--\didUpperPct{} percent in the structural window. I use that as
a sign, timing, and scale check, not as the main estimate. The central object
is a structural decomposition, under explicit auction assumptions, of how
price formation moves when the policy removes non-SMEs and reshuffles the
active SME pool.

The auction format makes the decomposition feasible. BEC's main procurement
format is the electronic \emph{Preg\~ao}, an English-reverse auction in which
firms submit successively lower offers until one remains. Under independent
private values, losing bidders' drop-out prices reveal their costs
\citep{haile2003}. That is the identifying advantage of the setting and also
its main structural restriction. I use drop-outs to recover type-specific cost
distributions for SMEs and non-SMEs, correct for auction-level scale
heterogeneity in the spirit of \citet{krasnokutskaya2011}, and simulate
counterfactual auctions under observed equilibrium entry. The exercise is not
a full entry equilibrium and not a reduced-form treatment-effect estimate:
observed pre- and post-policy entry rates are taken as primitives, and the
post-policy SME cost distribution is estimated from post-policy exits rather
than derived from a formal selection model.

The decomposition uses three counterfactual prices. Let $S_1$ be the open
auction with the pre-policy bidder pool, $S_2$ the SME-only auction holding
the pre-policy SME pool fixed, and $S_3$ the SME-only auction with the
observed post-policy SME pool. Then
\[
    p_{S_3} - p_{S_1}
    =
    \underbrace{p_{S_2} - p_{S_1}}_{\text{lost competitive discipline}}
    +
    \underbrace{p_{S_3} - p_{S_2}}_{\text{protected-pool offset}}.
\]
The first term measures the cost of removing non-SMEs from the price-forming
pool. The second measures the offset created by the post-policy protected
pool, which combines additional SME participation with changes in the
composition of active SME bidders. The decomposition turns a set-aside price
effect into a design question: how much of the cost comes from replacing
competition by eligibility, and how much is recovered by the protected pool
that the policy brings in?

The main result is that the protected pool responds, but exclusion dominates.
In non-pharmaceutical supplies, removing non-SMEs while holding the pre-policy
SME pool fixed raises simulated prices by \bneEffectIntensNp{} of the
reference price; the observed post-policy SME pool offsets that shock by
\bneEffectEntryNp{}, leaving a net effect of \bneEffectTotalNp{}, with the
exclusion channel accounting for \bneShareIntensNp{} percent of the two
components in absolute magnitude. The protected pool attracts and reshuffles
SMEs, but it does not recreate the competitive discipline supplied by
excluded non-SMEs. Pharmaceutical procurement shows the same qualitative
pattern at larger magnitudes, with composition turnover that makes the
welfare ranking conditional rather than headline.

This price-formation result changes the policy comparison. A full set-aside
supports SMEs by removing their rivals; a price preference does so while
keeping non-SMEs in the auction. As a static design benchmark, I use a 10
percent SME price preference under fixed entry and cost primitives. In thick,
standardized non-pharmaceutical markets, the preference has near-zero price
and welfare cost while the full set-aside generates a substantial loss. The
preference is not a free version of the set-aside; it delivers smaller
redistribution. The comparison is therefore a frontier, not a ranking: a
planner should choose the set-aside only when the additional SME surplus is
worth the allocative and fiscal cost of removing non-SMEs.

The strongest policy result is in thick standardized markets. In thinner and
more heterogeneous pharmaceutical markets, the ranking depends on how the
post-policy SME pool is modeled. The main specification estimates the
post-policy SME cost distribution from post-policy exits; a strict invariance
benchmark instead holds the SME cost distribution fixed across the cutoff.
The non-pharmaceutical ranking is stable across these readings; the
pharmaceutical ranking is not. I treat the gap as a scope condition rather
than a second headline: bidder exclusion is hardest to defend where the
protected pool is thick enough for a preference rule to tilt allocation
without removing competitive discipline, and most defensible, if at all,
where the protected pool is thin, the good is heterogeneous, and the planner
places high value on SME surplus.

The paper develops that frontier through layered analysis built around a
single engine. The engine is the decomposition: the set-aside's price effect
splits into a discipline-loss channel from removing non-SMEs and a
protected-pool offset from the SME response. Drop-out bids in the electronic
Preg\~ao give the channels empirical content without invoking bid-function
inversion \citep{krasnokutskaya2011,krasnokutskayaseim2011}, and a planner
problem then prices the channels against transfers to SMEs, allocative waste,
and the marginal cost of public funds \citep{saezstantcheva2016}. The lesson
is direct: when the bidders a set-aside removes are the ones forming the
price, the rule trades competitive discipline for an incomplete protected-pool
response and earns the difference as a welfare loss. Full set-asides remain
defensible, but the defense requires the planner to assign explicit weight to
SME surplus, market access, or dynamic capacity gains beyond what the
protected pool itself delivers. The natural benchmark for SME procurement
support is not no support, but support that keeps rival bidders inside the
auction.
